% ---------------------------------------------------------------------------
% Chapter: Introduction
% ---------------------------------------------------------------------------
\chapter{Introduction}\label{chap:intro}
% Long monospaced identifiers leave little room for line breaking; allow some
% extra inter-word stretch in this chapter (the group keeps it local).
\begingroup\setlength{\emergencystretch}{3em}% chapter-local

My Tool Box is a Python package for research on networked multi-agent control
and multi-agent reinforcement learning (MARL). It is distributed as a single
installable project, \pkg{my-tool-box}, that provides two import packages:

\begin{description}[leftmargin=1.6cm,style=nextline]
  \item[\pkg{env\_lib}] Multi-agent environments that follow the Gymnasium
    API: eight environment families registered under 18 ids, among them
    coupled oscillators, consensus and formation control, power-grid
    frequency control, vehicle platoons and multi-robot target tracking.
    Around the environments it provides natively batched vector
    environments, a searchable catalogue and the \code{env-lib} command line,
    adapters for other interfaces, a classical baseline controller for every
    environment, evaluation tools and shared graph utilities.
  \item[\pkg{toolkit}] Research utilities that do not depend on the
    environments: \pkg{plotkit} creates publication-quality figures,
    \pkg{neural\_toolkit} provides PyTorch network building blocks and
    tabular reinforcement-learning tools, and \pkg{parakit} manages
    \pkg{argparse} parameters, with an optional graphical editor.
\end{description}

This manual describes version 1.1 of the package. It explains how to install
it, documents every environment and toolkit module and the workflow layer that
connects them to experiments, collects practical recipes and troubleshooting
advice, and ends with an API reference and a migration guide for code written
against earlier versions.

\section{What the package is built for}\label{sec:intro-audience}

The package is built for research on control and learning in networks of
interacting agents. Its environments share five characteristics.

\begin{itemize}
  \item \emph{Networked agents.} Agents interact through an explicit graph:
    a communication topology, the lines of a transmission network or the
    radio links of a platoon. Each agent observes its own state and its
    neighbourhood, so the environments pose decentralised problems in which
    the graph matters.
  \item \emph{Continuous physical dynamics.} Most environments have
    continuous states and actions in physical units, with dynamics that come
    from established models: phase oscillators, single and double
    integrators, the swing equations of a power system, the longitudinal
    dynamics of vehicles with actuator lag, and robot kinematics with
    estimation. Models, observations, rewards and parameters are documented
    precisely enough to be cited in a paper.
  \item \emph{Many agents.} Team sizes are constructor arguments, and
    dynamics, observations and rewards are computed with array expressions
    rather than loops over agents, so an environment with a hundred agents
    remains practical.
  \item \emph{Large-batch simulation.} The networked continuous-control
    environments (Kuramoto, Consensus and Formation, PowerGrid, Platoon) have
    native vector implementations that advance thousands of copies with one
    set of array operations per step, which suits on-policy learning with
    large batches, Monte-Carlo evaluation and parameter sweeps.
  \item \emph{Reproducibility.} All randomness is seeded, benchmark systems
    such as a power network or a heterogeneous platoon are fixed by their own
    seeds, a single environment and copy~0 of its vector environment share
    their code, and every environment ships a classical baseline controller
    that provides a reference result.
\end{itemize}

Two discrete networked benchmarks (message relaying and wireless access) and a
cooperative physics game complement the continuous-control environments. The
environments do not depend on any particular learning library; any algorithm
that can drive a Gymnasium environment can use them, and the adapters of
Chapter~\ref{chap:workflow} provide flat single-agent and per-agent dictionary
views of the same environments.

\section{What the package provides}\label{sec:intro-provides}

Table~\ref{tab:intro-envs} lists the environments. Each one is available as a
class in \pkg{env\_lib} and under one or more registered Gymnasium ids; the
version suffix of an id selects a preset configuration (for example
\envid{WirelessComm-v1} is the $4\times 4$ grid), not a newer revision.
\code{env\_lib.list\_envs()} returns all 18 ids, and \code{env-lib list} prints
them as a table. Figure~\ref{fig:intro-gallery} shows the built-in
visualisations.

\begin{table}[htbp]
  \centering
  \small
  \caption{Environments provided by \pkg{env\_lib}. ``Batch'' marks the
    environments with a native vector implementation
    (Chapter~\ref{chap:workflow}); the last column names the optional
    dependency group (Chapter~\ref{chap:install}) required in addition to the
    core installation.}
  \label{tab:intro-envs}
  \begin{tabularx}{\textwidth}{@{}>{\raggedright\arraybackslash}p{4.8cm}>{\raggedright\arraybackslash}X>{\raggedright\arraybackslash}p{0.7cm}cc@{}}
    \toprule
    Registered ids & Task & Ch. & Batch & Extra \\
    \midrule
    \multicolumn{5}{@{}l}{\emph{Continuous states and actions}} \\
    \envid{KuramotoOscillator-*} (4 ids) &
      Drive a network of coupled phase oscillators to synchronisation (NumPy). &
      \ref{chap:kuramoto} & yes & -- \\
    \envid{KuramotoOscillatorTorch-*} (5 ids) &
      The same model on PyTorch tensors, with many systems simulated in a batch. &
      \ref{chap:kuramoto} & -- & \code{torch} \\
    \envid{Consensus-v0}, \envid{Formation-v0} &
      Rendezvous or form a shape over a communication graph. &
      \ref{chap:consensus} & yes & -- \\
    \envid{PowerGrid-v0} &
      Arrest and damp frequency deviations of a power network after load
      changes. &
      \ref{chap:power_grid} & yes & -- \\
    \envid{Platoon-v0} &
      Keep a vehicle platoon at safe, string-stable spacing behind a leader. &
      \ref{chap:platoon} & yes & -- \\
    \envid{AJLATT-v0} &
      Robots localise themselves and jointly track a moving target. &
      \ref{chap:ajlatt} & -- & -- \\
    \envid{Pistonball-v0} &
      Pistons cooperate to roll a ball to the goal (2-D physics; continuous or
      discrete actions). &
      \ref{chap:pistonball} & -- & \code{pistonball} \\
    \midrule
    \multicolumn{5}{@{}l}{\emph{Discrete states and actions}} \\
    \envid{LineMsg-v0} &
      Relay a message along a line of agents with unreliable links. &
      \ref{chap:networked} & -- & -- \\
    \envid{WirelessComm-v0}, \envid{-v1} &
      Share wireless access points on a grid without collisions. &
      \ref{chap:networked} & -- & -- \\
    \bottomrule
  \end{tabularx}
\end{table}

\begin{figure}[htbp]
  \centering
  \newcommand{\galleryimg}[2]{%
    \begin{minipage}[t]{0.32\textwidth}\centering
      \includegraphics[width=\linewidth,height=3.1cm,keepaspectratio]{#1}\\[2pt]
      {\footnotesize #2}
    \end{minipage}}%
  \galleryimg{kuramoto.png}{(a) Kuramoto oscillators}\hfill
  \galleryimg{formation.png}{(b) Formation control}\hfill
  \galleryimg{power_grid.png}{(c) PowerGrid}\\[8pt]
  \galleryimg{platoon.png}{(d) Platoon}\hfill
  \galleryimg{ajlatt.png}{(e) AJLATT target tracking}\hfill
  \galleryimg{pistonball.png}{(f) Pistonball}\\[8pt]
  \galleryimg{linemsg.png}{(g) LineMsg}\hspace{0.02\textwidth}%
  \galleryimg{wireless.png}{(h) WirelessComm}
  \caption{Frames rendered in \code{"rgb\_array"} mode by the environments of
    \pkg{env\_lib}. Each renderer is a dashboard that combines the scene with
    the quantities that matter for the task, such as the order parameter of the
    oscillators, the frequency nadir of a power grid, the spacing errors along
    a platoon, per-agent outcomes or covariance traces.}
  \label{fig:intro-gallery}
\end{figure}

Beyond the environments, \pkg{env\_lib} provides the workflow layer of
Chapter~\ref{chap:workflow}: \code{env\_lib.make\_vec} for batched simulation,
\code{env\_lib.catalog} and \code{env\_lib.describe} to discover environments,
the \code{env-lib} command line to list, run, record, evaluate and benchmark
them without writing code, wrappers that flatten the joint spaces or expose
per-agent dictionaries, \code{env\_lib.baseline\_policy} for the baseline
controllers, and \code{env\_lib.evaluate} and \code{env\_lib.rollout} for
evaluation and trajectory recording. The graph utilities of
\pkg{env\_lib.utils.graphs} (Section~\ref{sec:envlib-graphs}) build
topologies and compute Laplacians, algebraic connectivity and $k$-hop
neighbourhoods.

The \pkg{toolkit} package contains three independent subpackages:

\begin{itemize}
  \item \pkg{toolkit.plotkit} (Chapter~\ref{chap:plotkit}): learning curves
    with confidence bands and smoothing, grouped bar charts, histograms,
    annotated heatmaps, style presets, colour-blind-safe palettes and
    multi-format export. It depends on matplotlib only.
  \item \pkg{toolkit.neural\_toolkit} (Chapter~\ref{chap:neural}): policy,
    value and Q-networks with MLP, convolutional, recurrent and transformer
    trunks, encoders and decoders (including a variational autoencoder),
    factories, network utilities and tabular reinforcement-learning tools.
    It requires PyTorch.
  \item \pkg{toolkit.parakit} (Chapter~\ref{chap:toolkit}): saving, loading,
    validating and applying the defaults of an \code{argparse.ArgumentParser},
    and a Tk window for editing them.
\end{itemize}

\section{Design principles}\label{sec:intro-principles}

\paragraph{One Gymnasium contract.}
Every environment subclasses \code{gymnasium.Env}. The method \code{reset}
accepts a \code{seed} and an \code{options} dictionary as keyword arguments
and returns the pair \code{(observation, info)}. The method \code{step}
returns the observation, the reward, the flags \code{terminated} and
\code{truncated}, and the \code{info} dictionary, where \code{terminated}
signals a task-defined end (success or failure) and \code{truncated} the
episode limit. The whole team is exposed through \emph{joint} spaces whose
leading axis indexes the agents: observations are \code{float32} arrays of
shape \code{(n\_agents, obs\_dim)} contained in \code{observation\_space}, and
the action covers all agents at once. \code{step} returns the team reward, and
every environment reports the per-agent rewards in
\code{info["agent\_rewards"]}. Each environment enforces its own episode limit
through a constructor argument, and \code{env\_lib.make(...,
max\_episode\_steps=N)} sets that argument instead of adding a second limit.
Chapter~\ref{chap:envlib} describes these conventions in detail.

\paragraph{Batch-first kernels shared by single and vector environments.}
The dynamics of the environments with a native vector implementation are
written once, as array operations on a leading batch dimension $B$. The single
environment runs them with $B = 1$ and the native vector environment with
$B$ = \code{num\_envs}, so both follow exactly the same model, and copy~0 of a
vector environment reproduces a single environment started from the same state
(the test suite checks this). Their steps contain no Python loop over agents or
copies. The Kuramoto model additionally has a PyTorch backend that simulates many systems
on the CPU or a GPU. Section~\ref{sec:examples-benchmarks} and
Chapter~\ref{chap:workflow} report measured step times and throughput.

\paragraph{The interaction graph as data.}
Who observes, communicates with or is coupled to whom is explicit data rather
than a hidden detail of the code. The networked continuous-control
environments expose the graph as a matrix (\code{adjacency} of Consensus,
PowerGrid and Platoon, \code{topology\_matrix} and \code{coupling\_matrix} of
the Kuramoto environments), together with the Laplacian or the line
susceptances where they matter, and document the layout of the observation in
an \code{observation\_layout} property. The shared functions of
\pkg{env\_lib.utils.graphs} build standard topologies (PowerGrid draws its
transmission network with them) and compute Laplacians, algebraic
connectivity and $k$-hop neighbourhoods of any of these matrices.

\paragraph{Baselines as functions of the observation.}
Every environment ships a classical controller, decentralised wherever the
environment has several agents: Laplacian consensus, droop control,
cooperative adaptive cruise control and others
(Section~\ref{sec:workflow-baselines}). Each controller is a pure function of
the observation that accepts any leading batch shape, so the same object
drives a single environment and all copies of a vector environment, and a
decentralised controller computes the action of an agent from that agent's
row of the observation only.
\code{env\_lib.baseline\_policy(env)} returns the controller for any
environment, which makes a reference result one line of code.

\paragraph{Light core, lazy extras.}
\code{import env\_lib} and \code{import toolkit} load only the core
dependencies (NumPy, SciPy, matplotlib, Gymnasium and PyYAML). Environment
classes, renderers and toolkit subpackages are imported on first use, and
PyTorch, pygame, pymunk, imageio and Tk are imported only when a feature that
needs them is used; a missing dependency produces an error message that names
the package extra to install. Other packages are optional integrations, never
requirements: the parallel adapter of Section~\ref{sec:workflow-wrappers},
for example, works without the library whose interface it implements.

\paragraph{Reproducibility.}
All randomness of an environment comes from its own random generator
(\code{self.np\_random}), which is seeded by \code{reset(seed=...)}. No library
code seeds or draws from the global NumPy, Python or PyTorch random state, and
plotting functions never modify the global matplotlib configuration unless
asked to. Section~\ref{sec:trouble-seeding} lists everything that must be
seeded for an experiment to be repeatable.

\paragraph{Headless rendering.}
Rendering is optional and never performed unless requested. In
\code{"rgb\_array"} mode figures are drawn off-screen and returned as
\code{(H, W, 3)} \code{uint8} arrays, so episodes can be recorded to GIF or MP4
on machines without a display, for example on a compute cluster or in
continuous integration. In \code{"human"} mode an environment renders itself
at the end of every \code{reset()} and \code{step()}, as Gymnasium prescribes.

\section{Package layout}\label{sec:intro-layout}

The repository uses the standard \code{src} layout:

\begin{textcode}
My_Tool_Box/
  pyproject.toml           packaging metadata, optional extras, tool settings
  CHANGELOG.md             release notes and migration notes
  LICENSE                  MIT licence
  src/
    env_lib/               environments (import name env_lib)
      registration.py      ids and metadata, make(), make_vec(), get_spec()
      catalog.py           catalog(), describe()
      baselines.py         baseline controllers, baseline_policy()
      __main__.py          the env-lib command line
      errors.py            ResetNeededError
      wrappers/            FlattenJointSpaces, TeamReward, ParallelEnvAdapter
      utils/               rendering, recording, graphs, vector base class,
                           evaluate() and rollout()
      kos_env/             Kuramoto oscillators (NumPy, PyTorch, vector)
      consensus_env/       Consensus and Formation (single and vector)
      power_grid_env/      PowerGrid (single and vector)
      platoon_env/         Platoon (single and vector)
      ajlatt_env/          AJLATT (configuration, estimation, maps, rendering)
      pistonball_env/      Pistonball
      linemsg_env/         LineMsg
      wireless_comm_env/   WirelessComm
    toolkit/               research utilities (import name toolkit)
      plotkit/             plotting
      neural_toolkit/      PyTorch networks and tabular tools
      parakit/             argparse parameter management
  examples/                runnable, documented example scripts
  benchmarks/              performance measurements
  tests/                   pytest suite
  docs/manual/             this manual (LaTeX sources)
\end{textcode}

A first program creates an environment through the registry, runs its
baseline controller for one episode, and then evaluates the same controller on
256 copies of the environment simulated as one batch:

\begin{pycode}
import env_lib

env = env_lib.make("PowerGrid-v0")              # 16 buses, one agent per bus
policy = env_lib.baseline_policy(env)           # decentralised droop control
obs, info = env.reset(seed=0)                   # obs: (16, 10) float32
while True:
    obs, reward, terminated, truncated, info = env.step(policy(obs))
    if terminated or truncated:
        break
print(info["peak_abs_omega"], info["agent_rewards"].shape)   # rad/s, (16,)

envs = env_lib.make_vec("PowerGrid-v0", num_envs=256)       # native batch
result = env_lib.evaluate(envs, env_lib.baseline_policy(envs),
                          n_episodes=256, seed=0)
print(result.mean_return, result.ci95)          # mean return, 95% CI half-width
\end{pycode}

The same can be done without code: \code{env-lib run PowerGrid-v0} runs one
episode of the baseline and \code{env-lib evaluate PowerGrid-v0 -{}-num-envs 64}
compares it with random actions (Section~\ref{sec:workflow-cli}).

\section{Conventions used in this manual}\label{sec:intro-conventions}

Python identifiers are set in \code{monospace}, package and module names in
\pkg{sans-serif}, and registered environment ids as \envid{Consensus-v0}.
Shell commands are meant to be run from the repository root in a POSIX shell;
on Windows, replace \code{export VAR=value} by \code{set VAR=value}
(\code{cmd}) or \code{\$env:VAR="value"} (PowerShell). Code examples use the
public API of version 1.1 and were executed with it. Array shapes are written
as tuples, for example \code{(n\_agents, obs\_dim)}; \code{B} denotes a batch
size (the number of copies of a vector environment) and \code{T} a sequence
length.

\section{How this manual is organised}\label{sec:intro-organisation}

\begin{itemize}
  \item Chapter~\ref{chap:install} covers requirements, installation with
    optional extras, verification and building this manual.
  \item Part~I documents the environments: the common interface, the
    registry and the graph utilities (Chapter~\ref{chap:envlib}), the
    Kuramoto oscillators (Chapter~\ref{chap:kuramoto}), LineMsg and
    WirelessComm (Chapter~\ref{chap:networked}), Pistonball
    (Chapter~\ref{chap:pistonball}), consensus and formation control
    (Chapter~\ref{chap:consensus}), power-grid frequency control
    (Chapter~\ref{chap:power_grid}), vehicle platoons
    (Chapter~\ref{chap:platoon}) and AJLATT (Chapter~\ref{chap:ajlatt}). It
    closes with the workflow layer: vectorised simulation, the catalogue and
    the command line, adapters, baseline controllers and evaluation
    (Chapter~\ref{chap:workflow}).
  \item Part~II documents the toolkit: its structure and \pkg{parakit}
    (Chapter~\ref{chap:toolkit}), \pkg{plotkit} (Chapter~\ref{chap:plotkit})
    and \pkg{neural\_toolkit} (Chapter~\ref{chap:neural}).
  \item Part~III is practical: the example scripts, recording, benchmarks and
    the contributor workflow (Chapter~\ref{chap:examples}), and
    troubleshooting (Chapter~\ref{chap:troubleshooting}).
  \item Appendix~\ref{app:api} is a compact API reference and
    Appendix~\ref{app:migration} explains how to migrate code written for
    version 1.0 and for the versions before it.
\end{itemize}

\section{Licence and citation}\label{sec:intro-licence}

My Tool Box is free software distributed under the MIT licence; the full
text is in the file \code{LICENSE} at the root of the repository. The source
code is hosted at \url{https://github.com/EricDMW/My_Tool_Box}, where bugs
and feature requests can be reported as issues.

If the package contributes to published work, please cite it as follows:

\begin{textcode}
@software{wang2026mytoolbox,
  author  = {Wang, Dongming},
  title   = {My Tool Box: Multi-Agent Reinforcement Learning
             Environments and Research Utilities},
  version = {1.1.0},
  year    = {2026},
  url     = {https://github.com/EricDMW/My_Tool_Box}
}
\end{textcode}

When results depend on a particular environment, please also cite the
original description of the underlying model or benchmark (for example the
Kuramoto model of coupled oscillators), as given in the corresponding
environment chapter.

\par\endgroup
